\chapter{Introduction}
\label{sec:introduction}

% TODO(bibliography): citations below use \citep{} keys not yet in
% bibliographies/references.bib, so they render as [?] until the bibliography pass.
% Keys used here:
%   robodopamine -> RoboDopamine
%   roboreward   -> RoboReward
%   lrm          -> "Large Reward Models: Generalizable Online Robot Reward
%                    Generation with Vision-Language Models"
%   robometer    -> Robometer (Liang et al.)        [the base model we build on]
%   demo2reward  -> demo2reward (Gumbsch et al., 2026) [overconfidence claim]
%   brown2020    -> Brown et al., 2020, "Language Models are Few-Shot Learners"
%   dong2022     -> Dong et al., 2022, "A Survey on In-Context Learning"
% The Abstract still uses plain-text "(Liang et al.)"/"(Gumbsch et al., 2026)";
% convert both to \citep in the same pass.

Reinforcement learning (RL) offers a principled framework for teaching robots to act through trial and error, but it rests on a deceptively difficult prerequisite: a reward function that specifies what the robot should achieve. For robotic manipulation, designing this signal by hand is notoriously laborious and brittle. Dense, shaped rewards demand task-specific engineering and privileged state information --- object poses, contact flags, goal distances --- that is rarely available outside simulation, while sparse success rewards are easy to specify but hard to optimize. Worse, a reward crafted for one task seldom transfers to the next, so every new behavior reopens the same engineering problem. In practice the reward, rather than the policy optimizer, is often the true bottleneck: given a faithful reward the same RL loop learns readily, and without one it fails regardless of the algorithm.

Vision-language models (VLMs) trained on internet-scale data have recently shown broad visual and semantic understanding that transfers across domains with little or no task-specific tuning. This raises the prospect of a general, reusable reward: a model that simply watches an episode and judges whether --- and how much --- the task was accomplished, replacing per-task reward engineering with a single pretrained judge. It also poses the central question of this thesis: \textit{can state-of-the-art VLMs provide meaningful, generalizable rewards and resolve the reward-definition bottleneck in robot RL?}

Several recent works pursue this idea, training VLMs to predict rewards for robot behavior \citep{robodopamine, roboreward, lrm}. The most competitive of these is Robometer \citep{robometer}, a large multi-task VLM reward model trained on a corpus of over a million robot trajectories and equipped with dedicated output heads for success and task-progress prediction. Robometer reports strong offline reward-quality metrics, yet the field remains young and largely unexplored. In particular, these models are validated almost entirely by offline metrics computed on held-out clips; whether they actually function as reward signals inside a downstream RL loop --- the setting they are ultimately meant for --- has received far less scrutiny.

We introduce \textbf{RoboRef}, a VLM reward model that builds on Robometer through three contributions. The first, and most central, concerns the training data. Robometer supervises its two most important heads --- success and task progress --- almost exclusively on \emph{successful} trajectories; failures reach the model only indirectly, as the rejected side of a preference comparison. A reward model that has never been densely trained on failure has little reason to recognize it, yet recognizing failure is exactly what an RL policy will stress-test. We therefore drop the preference head and supervise the success and progress heads directly on failure trajectories. Doing so requires what, to our knowledge, does not yet exist: a robot dataset in which failed attempts carry \emph{dense, per-frame} progress annotations. Annotating failures densely is the crux of the problem. For a successful trajectory, progress can be approximated by how far the demonstration has advanced; a failed attempt has no such monotonic reference --- it may stall, regress, or briefly approach success before collapsing. We obtain reliable labels from two complementary sources. In simulation, we read the simulator's privileged state at every step, producing exact per-frame progress for a controlled spectrum of failure modes --- from an arm that never reaches the object, to a grasp that slips moments before completion. For real-world failures, where no such ground truth exists, we design a dual pipeline: a vision-language model first describes each frame in isolation, and a language model then acts as a judge that maps those descriptions to a per-frame progress score. Decoupling perception from adjudication in this way proved substantially more reliable than asking a single model to score an entire video.

Our second contribution introduces in-context demonstrations to the reward model. In-context learning --- supplying examples in the prompt so that a model adapts to a task without any weight update --- is among the defining capabilities of large pretrained models \citep{brown2020}, and has been shown to improve few-shot generalization across language and multimodal tasks \citep{dong2022}. For a reward model, a demonstration supplies precisely what a textual task description cannot: a concrete, visual definition of what success looks like for this embodiment, scene, and camera viewpoint. We pair each trajectory under evaluation with a successful demonstration of the same task and, to our knowledge, are the first to do so for a robot VLM reward model. This grounds the model's judgment in an example rather than in a potentially ambiguous instruction, sharpening its notion of task completion.

Our third contribution revisits the training objective. VLM reward models have been observed to be overconfident, systematically over-predicting success \citep{demo2reward}. In reinforcement learning, the two ways of being wrong are far from symmetric. A false positive --- the reward declaring success when the task has actually failed --- creates a spurious goal that an optimizing policy will actively seek out and exploit, the failure mode known as reward hacking. A false negative --- withholding reward from genuine success or progress --- is comparatively harmless: it can slow learning by making the reward signal sparser, but it never hands the policy a shortcut to exploit. A reward model used for RL should therefore err toward caution. We introduce a targeted, principled modification of the standard reward objective that penalizes over-prediction more heavily than under-prediction, biasing the model toward conservative, harder-to-exploit predictions.

We evaluate RoboRef both offline and as the reward signal for downstream RL, across two simulated manipulation benchmarks. Our experiments expose a sharp disconnect: reward models that excel on standard offline metrics can nonetheless fail completely as RL rewards, because an optimizing policy learns to exploit their false positives. RoboRef, by contrast, trains policies on tasks where the off-the-shelf baseline never learns at all, and outperforms prior reward models in both the offline and the downstream regime.

In summary, this thesis makes the following contributions:
\begin{itemize}
    \item \textbf{A failure-annotated reward dataset.} To our knowledge the first robot dataset with dense, per-frame progress labels on \emph{failure} trajectories, built from simulation ground truth and a dual VLM-then-LLM labeling pipeline.
    \item \textbf{In-context demonstrations for reward modeling.} A novel addition relative to prior VLM reward models, grounding the model's notion of task success in a concrete example rather than a text instruction.
    \item \textbf{An asymmetric reward objective.} A targeted modification of the standard loss that addresses the documented overconfidence of VLM reward models by penalizing false positives, yielding more conservative and less exploitable rewards.
\end{itemize}
Along the way, we surface an observation with implications beyond our own model: the offline metrics reported by current VLM reward models do not predict their usefulness as downstream RL rewards --- a gap we examine directly.

These contributions and observations are organized around the following \textit{preliminary} research questions. Our main research question is:

\noindent\textit{To what extent can a fine-tuned vision-language model serve as a usable and generalizable reward signal for downstream robot reinforcement learning?}

\noindent which we approach through four sub-questions:
\begin{itemize}
    \item \textbf{RQ1.} Does supervising the reward model on curated failure trajectories with dense per-frame progress labels improve its ability to discriminate success from failure?
    \item \textbf{RQ2.} Are VLM reward models overconfident, and does an asymmetric loss that penalizes false positives yield more conservative, less exploitable rewards?
    \item \textbf{RQ3.} Do offline reward-quality metrics predict downstream RL performance?
    \item \textbf{RQ4.} Under what conditions can a fine-tuned VLM reward model train a policy where state-of-the-art baselines fail?
\end{itemize}
